\section{Open obligations}\label{sec:open-obligations}

The proved local core is confined to the formal Hamiltonian,
residue, factorization-current, scalar-reduced Moyal, and scalar-anomaly
claims stated above. The following statements are outside those
theorems.

\begin{enumerate}
  \item \emph{Proof ownership behind theorem-lane entries.}
  Each index lane requires a cited theorem, proposition, lemma, or
  computation whose hypotheses and conclusion match the lane. The index
  entry itself is not a proof-bearing theorem.

  \item \emph{Sign and convention appendix.}
  A convention block must record the BV degrees, Ext shifts, CE/PV
  shifts, current-density pairings, interval-orientation sign, scalar
  quotient, and Capelli/Weyl normalization. It must distinguish
  convention signs from computed signs.

  \item \emph{Matlis/local-cohomology source boundary.}
  Appendix~\ref{sec:app-matlis-principal-parts} contains the topology on
  \(\mathfrak h\), the continuous dual calculation, the residue action,
  the local-cohomology uniqueness statement, and the
  no-polynomial-descendant theorem.  It also contains the limited
  Fourier--Rees label bridge: the bridge is a Weyl-module
  Fourier-normalization and associated-graded label comparison through
  the Fourier delta Rees lattice, not an \(\mathfrak h\)-module
  identification and not an equality with the full Cech lattice before
  \(\hbar\) is inverted.  The remaining obligation is any theorem that
  changes the Hamiltonian action or the boundary target enough to evade
  the same-action local-nilpotence obstruction.

  \item \emph{Radial-parts source matching.}
  Appendix~\ref{sec:app-radial-parts-moyal} fixes the Weyl algebra
  conventions, Capelli triangular signs, scalar \(U(1)\) quotient, and
  relation to the stated finite exact ranges.  The finite-\(N\) all-\(N\)
  radial-parts identity remains conditional on matching the
  Harish-Chandra--Wallach--Levasseur--Stafford input to this quantum
  Hamiltonian reduction.  The finite exact check is independent finite-range
  evidence, not a substitute for that source match.

  \item \emph{Unreduced cotangent factorization centre.}
  An unreduced open BV model whose cotangent sector carries
  principal-part currents before de Rham-current contraction remains open.
  The corresponding centrality homotopies against arbitrary open
  observables are also open.  The obstruction is visible before any
  analytic graph integral is chosen.  In the grading used here,
  \(|\phi_i|=0\), \(\psi=A^*\) is odd of BV degree \(-1\),
  a Hamiltonian generator \(H_f\) has degree \(0\), a shifted cotangent
  generator \(\eta_\rho=\rho[1]\) has degree \(-1\), and the CE
  coordinates have \(|c_f|=1\), \(|u_\rho|=0\).  The polynomial
  one-\(\psi\) descendants cannot serve as unreduced cotangent
  representatives: for the linear Hamiltonians,
  \[
    \{O_{1,0},\widetilde\Psi_{a,b}\}=b\,\widetilde\Psi_{a,b-1},
    \qquad
    \{O_{0,1},\widetilde\Psi_{a,b}\}
      =-a\,\widetilde\Psi_{a-1,b},
  \]
  while the principal-part coadjoint action is
  \[
    z_1\cdot\rho_{a,b}=-(b+1)\rho_{a,b+1},
    \qquad
    z_2\cdot\rho_{a,b}=(a+1)\rho_{a+1,b}.
  \]
  The index shifts go in opposite directions.  Hence the primitive
  one-\(\psi\) projection is only the indecomposable \(P_0\) shadow of
  this problem; it is not a centrality homotopy and not an unreduced
  factorization-centre morphism.

  \item \emph{One-antifield quantum lift.}
  The extension of the degree-zero Moyal/Weyl comparison to the
  descendant sector is open at the unreduced open-BV and
  factorization-centre level. The finite nilpotent source
  \(\mathfrak m^3/\mathfrak m^{N+1}\) is not stable under the Moyal
  bracket for \(N\geq4\), so the lift requires a different
  Moyal-flat source or a weighted coefficient construction with its own
  hypotheses.  There is, however, a target-side reduced complex:
  Construction~\ref{constr:quantum-principal-part-descendant-target}
  defines the distributional principal-part target with full Moyal
  coadjoint action
  \[
    \{B_f(a),\Theta_\rho(B)\}_\hbar
      =
      \Theta_{\operatorname{ad}^{\vee}_{f,\hbar}\rho}(aB),
    \qquad
    \{\Theta_\rho(B),\Theta_\sigma(C)\}_\hbar=0.
  \]
  This target does not come from the primitive polynomial one-\(\psi\)
  classes.  Proposition~\ref{prop:quantum-descendant-label-obstruction}
  computes
  \[
    \operatorname{ad}^{\vee}_{z_1^4,\hbar}\delta_{1,1}
      =-24\hbar^2\delta_{0,4}+\cdots,
  \]
  whereas the PBW one-\(\psi\) action sends
  \(\widetilde\Psi_{1,1}\) to \(4\widetilde\Psi_{4,0}\).  Thus the
  obstruction survives after primitive projection and is not a
  multi-trace artifact.  Theorem~\ref{thm:no-hbar-primitive-descendant-intertwiner}
  strengthens this: no \(\hbar\)-adic deformation whose leading term is
  \(\widetilde\Psi_{a,b}\mapsto\rho_{a,b}\) and whose source linear
  Hamiltonians have the PBW leading action can intertwine the Moyal
  coadjoint action.  The remaining obligation is therefore not to add
  higher corrections to the old label map, but to construct the
  unreduced boundary model, centrality homotopies, and analytic
  Costello/QME lift landing in the principal-part target.

  \item \emph{Regulator independence and the strict endpoint.}
  The weighted Tate regulator is independent of the chosen admissible
  degree weight inside the finite-window Mittag-Leffler sector, by
  Theorem~\ref{thm:wt-regulator-independence-admissible}.  This is not
  a theorem that the original product/direct-sum Tate pair carries the
  same closed-open graph theory: in that strict topology the Costello
  coefficient Casimir is absent by
  Lemma~\ref{lem:tate-casimir-obstruction}.  The stronger strict
  endpoint is ruled out, under the Costello coefficient-kernel and
  finite-window compatibility hypotheses, by
  Theorem~\ref{thm:strict-unweighted-no-go}.  Tail-sensitive observables
  in the weighted product completion remain outside the proved
  regulator-admissible sector, and ordinary cohomology of the weighted
  product is not identified with the strict direct-sum dual.

  \item \emph{Mixed brane-defect QME.}
  The quantum master equation for the brane-defect Costello graph model,
  with the required boundary kernel and scalar-anomaly normalization, is
  open. The Hadamard/Mittag-Leffler theorem supplies the analytic
  reduction, not the vanishing of the obstruction class.  The residual
  class is the brane-defect local functional
  \[
    \operatorname{Ob}\in
    H^1\bigl(\mathcal O_{\mathrm{loc}}(\mathcal E_w),
      Q+\{I_0,-\}\bigr),
  \]
  together with the scalar contact normalization.  Classically, before
  scalar reduction,
  \[
    \{I_\partial(a,f),I_\partial(b,g)\}_{\mathrm{open}}
    =I_\partial(ab,\{f,g\}_{\bar A})
      +N\,\omega(f,g)\int_\R a(t)b(t)\,dt .
  \]
  Quantum mechanically the same boundary contact appears as
  \[
    \operatorname{Tr}(A\,XY)-\operatorname{Tr}(A\,YX)
      =\hbar N\,\operatorname{Tr}(A)
      \quad\bmod\ \mathcal W_N\widehat\mu(\mathfrak{gl}_N).
  \]
  Here \(X,Y\) are the Weyl generators corresponding to
  \(\phi_1,\phi_2\).  Known candidate QME lifts include cancellation by a
  closed exchange graph and counterterm, removal by an explicitly
  stated supertrace or central-character normalization, or absorption
  into renormalized boundary generators.  Exhaustiveness would require
  a computation of the relevant degree-one obstruction complex.

  \item \emph{All-order Costello graph realization.}
  An all-order renormalized BV graph model whose boundary product is the
  Weyl/Moyal product in the normalization used here remains open. The
  first and third graph normalizations remain computations to stated
  order.

  \item \emph{Rees/Fourier bridge between \(\Psi\) and \(\rho\).}
  The constructed Fourier--Rees comparison relates the PBW polynomial
  labels and the principal-part labels only after Fourier twist and
  degreewise Rees normalization:
  \[
    z_1^az_2^b\mapsto
    \partial_1^a\partial_2^be_0
    \mapsto
    (-1)^{a+b}\hbar^{a+b}a!b!\rho_{a,b}.
  \]
  Before \(\hbar\) is inverted this is the Fourier delta Rees lattice
  \(\mathcal D_\hbar/(z_1,z_2)\), not the full Cech lattice
  \(\bigoplus\C[[\hbar]]\rho_{a,b}\).  It does not intertwine the
  original Hamiltonian action: Fourier sends the Rees-normalized
  operators \(\hbar X_{z_1},\hbar X_{z_2}\) to multiplication
  operators, while the Matlis cotangent module uses the original
  coadjoint vector fields on principal parts.  A same-action bridge
  from \(\Psi_{a,b}\) to \(\rho_{a,b}\) is ruled out on the generic
  fibre by local nilpotence unless the Hamiltonian action or the
  boundary target is changed.  Theorem~\ref{thm:no-hbar-primitive-descendant-intertwiner}
  records the corresponding \(\hbar\)-adic no-go for deformations of
  the classical primitive projection.

  \item \emph{Claim-status vocabulary.}
  Theorem boxes, the claim-strength ledger, outlook material, and
  theorem-lane entries require the same status vocabulary:
  proved in finite algebraic model, proved degreewise stable, proved in
  weighted Tate coefficient model, computed to stated order, conditional
  on named QME/radial-parts/unweighted-regulator theorem, conjectural,
  open, or not asserted.

  \item \emph{Global and cross-volume transfer.}
  Any compact Calabi--Yau, global BCOV, Vol III chiral-homology,
  \(\Phi_d\), Igusa/BKM/Borcherds, \(K3\times E\), or
  higher-dimensional derived-centre consequence requires a separate
  matched-conventions theorem containing the datum of
  Definition~\ref{defn:matched-conventions-theorem-datum}.  The local
  Hamiltonian BF/Moyal calculation supplies no such transfer by itself,
  and may be cited by sibling volumes only as a local \(d=2\) convention
  check unless they supply their own proof.  The templates in
  Remark~\ref{rmk:matched-conventions-templates} are obligations, not
  theorems; citing the template or the no-transfer lemma does not create
  a comparison morphism.

  \item \emph{Rendered-page, table, and figure verification.}
  The rendered title and abstract pages, claim-strength ledger, local
  dictionary, principles, theorem lanes, open obligations, and first- and
  third-order figures require visual verification. The figure labels and
  captions must state the computed order and hypothesis they support. A
  successful TeX run by itself does not verify these pages.

  \item \emph{Notation and macro hygiene.}
  Rendered and support notation files should contain no stale macros,
  duplicated definitions, informal text, or dead notation. A symbol
  should appear in the support files only when it is used with the same
  meaning in the manuscript or retained for a declared local convention.

  \item \emph{Primary-source anchoring.}
  For each use of Witten centrality, Costello graph/QME formalism, BCOV,
  Loday--Quillen--Tsygan, Procesi--Razmyslov, Matlis duality, and
  radial parts, record the source theorem, the local convention
  translation, what is used, and what is not used. A name citation alone
  does not transfer a theorem into the local model.
\end{enumerate}
